\section{Понятие вероятности}

\Def Пусть $A$ -- множество. Тогда $2^A \eqd \{ B: B \subset A \}$.

\Note $A \in 2^A, \varnothing \in 2^A$.

В принципе термины теории вероятности достаточно широко используются в бытовых ситуациях, в разговорах, в выпусках новостей и в прочих житейских вещах, поэтому, наверное, интуитивное понимание, что такое <<вероятность>> у нас есть.
Из этого давайте сформулируем формально, что мы под этим имееем в виду.

Что такое \textit{вероятность}? Базовое восприятие подсказывает, что это некоторая \textit{частота} появления какого-то интересующего нас \textit{события}.
Например, возьмем монетку и будем ее подкидывать. Как часто будет выпадать решка? Ну примерно в половине случаев она выпадет. В этом случае мы говорим, что вероятность $50/50$ или же математически $\frac{1}{2}$.
Что здесь является \textit{событием}? У монетки есть два события, которые нас интересуют -- выпадение орла либо решки. Это, так называемые, \textit{элементарные события}, т.е. которые, в рамках нашей модели, как бы неделимы.
То есть делаем вывод, что мы должны научится как-то мерить вероятность по крайней мере элементарных событий. Однако можно придумать более сложное событие, например, <<выпал орел или решка>>.
Ясно, что это будет в $100\%$ случаях или с вероятностью 1. Ну и также можно такой же трюк провернуть с союзом <<и>>. Было бы также интересно накладывать отрицания на события и тоже мерить их вероятность.

Теперь после этих рассуждений о том, чего мы хотим, мы готовы приступить к построению определений. Множество всех элементарных событий обозначают традиционно как $\Omega$ (в общем случае мы не накладываем никаких ограничений на него, это даже не факт что числа), а конкретное элементарное событие -- $\omega$.
И тогда построение комплексных событий посредством союзов <<и>>, <<или>>, <<не>> сводится к множественным операциям объединения, пересечения и дополнения.
Тогда вероятность можно воспринимать как фнукцию определенную на каком-то классе подножеств $\Omega$, то есть, формально, $\Pr: \mathcal{F} \to [0, 1]$. Где $\mathcal{F} \subset 2^\Omega$.
Однако, мы же хотим измерять вероятность объединений, пересечений и дополнений, т.е. функции вида $\Pr[A \cup B], \Pr[A \cap B], \Pr\left[A^C\right], \Pr\left[A^C \cap B\right]$, поэтому резонно наложить кое-какие ограничения на $\mathcal{F}$, чтобы \textit{результаты попадали в область определения вероятности}.

\Def Пусть $\Omega$ -- множество и $\mathcal{F} \subset 2^\Omega$. $\mathcal{F}$ будем называть \textit{$\sigma$-алгеброй} (сигма-алгеброй), если она удовлетворяет следующим аксиомам:
\begin{enumerate}
    \item $\Omega \in \mathcal{F}$.
    \item Пусть $\left\{ A_n \right\} \subset \mathcal{F}$. Тогда $\bigcup\limits_{n = 1}^\infty A_n \in \mathcal{F}$.
    \item Пусть $A \in \mathcal{F} \Rightarrow A^C \in \mathcal{F}$.
\end{enumerate}

\Exe Пусть $\mathcal{F}$ -- сигма-алгебра и $\left\{ A_n \right\} \subset \mathcal{F}$. Покажите, что $\bigcap\limits_{n = 1}^\infty A_n \in \mathcal{F}$.

Теперь логично функцию вероятности задавать на сигма-алгебре. Однако то что мы задали ее на естественной конструкции еще не означает, что мы достигли цели. Ведь в общем случае если для двух событий верно, что $A \cap B = \varnothing$, то из этого не следует, что $\Pr[A \cup B] = \Pr[A] + \Pr[B]$, а очень хотелось бы.
Более того, нам никто не гарантирует столь желаемого, что вероятность $\Omega$ равна 1.
Поэтому теперь введем ограничения на функцию.
\Def Пусть $\Omega$ -- множество и $\mathcal{F}$ -- сигма-алгебра над $\Omega$. Функцию $\Pr: \mathcal{F} \to [0, 1]$ будем называть \textit{вероятностью} (или вероятностной мерой), если она удовлетворяет следующим аксиомам:
\begin{enumerate}
    \item $\Pr[\Omega] = 1$
    \item $\Pr[\varnothing] = 0$
    \item Пусть $\left\{ A_n \right\} \subset \mathcal{F}: i \neq j \Rightarrow A_i \cap A_j = \varnothing$. Тогда:
    \begin{equation*}
        \Pr\left[ \bigcup\limits_{n = 1}^\infty A_n \right] = \sum\limits_{n = 1}^\infty \Pr [A_n]
    \end{equation*}
\end{enumerate}

\Def Пусть $\Omega$ -- множество, $\mathcal{F}$ -- сигма-алгебра над $\Omega$ и $\Pr: \mathcal{F} \to [0, 1]$ -- вероятность. Тогда тройка $(\Omega, \mathcal{F}, \Pr)$ называется \textit{вероятностным пространством}.

Следующие упражнения будут полезны, чтобы прочувствовать сущность предыдущих понятий.

\Exe Пусть $(\Omega, \mathcal{F}, \Pr)$ -- вероятностное пространство. Покажите, что $A \subset B \Rightarrow \Pr[A] \leq \Pr[B]$.

\Exe Постройте формально из нашего примера с монеткой вероятностное пространство. Покажите, что оно удовлетворяет всем аксиомам.

\Exe Пусть $(\Omega, \mathcal{F}, \Pr)$ -- вероятностное пространство и $A \in \mathcal{F}$. Покажите, что $\Pr\left[A^C\right] = 1 - \Pr[A]$.

\Exe Пусть $\Lambda$ -- непустое множество и $\left\{ \mathcal{F}_\lambda \right\}_{\lambda \in \Lambda}$ -- семейство $\sigma$-алгебр. Рассмотрим случаи:
\begin{enumerate}
    \item $\Lambda$ несчётно. 
    \item $\Lambda$ счётно.
    \item $\Lambda$ конечно.
\end{enumerate}
Можно ли хотя бы в одном случае сказать, что $\bigcup\limits_{\lambda \in \Lambda} \mathcal{F}_\lambda$ есть $\sigma$-алгебра?

\Exe Какие условия нужно наложить на $\left\{ \mathcal{F}_\lambda \right\}_{\lambda \in \Lambda}$ из предыдущего упражнения, чтобы всё-таки $\bigcup\limits_{\lambda \in \Lambda} \mathcal{F}_\lambda$ стало $\sigma$-алгеброй?

\section{Случайные величины}
\label{sec:random-variables}

В алгоритмах нас как правило интересуют числовые характеристики алгоритмов, такие как количество шагов. Рассмотрим отсортированный массив $A$ и напишем бинарного поиска элемента $t$.


\begin{algorithm}[H]
\caption{Бинарный поиск}
\begin{algorithmic}[1]

\Function{BIN\_SEARCH}{$A, t$}
    \State $l \gets 1$
    \State $r \gets |A|$

    \While{$l \leq r$}
        \State $m \gets \left\lfloor\frac{l + r}{2}\right\rfloor$
        \If{$A[m] = t$}
            \State \Return $m$
        \ElsIf{$A[m] < t$}
            \State $l \gets m + 1$
        \Else
            \State $r \gets m - 1$
        \EndIf
    \EndWhile

    \State \Return $-1$

\EndFunction
\end{algorithmic}
\end{algorithm}

Если нам приходит на вход какой-то случайный массив $A$, то количество шагов, которое выполнит бинарный поиск, логично, тоже будет случайным.
Пусть мы работаем на каком-то вероятностном пространстве $(\Omega, \mathcal{F}, \Pr)$. Чем характеризуется случайный массив? В рамках нашей модели он характеризуется каким-то случайным исходом $\omega \in \Omega$, в результате которого и получился такой массив (например, если за окном дождь, то он будет состоять из чисел Фибоначчи, а если солнце -- из нулей).
Но тогда получается, что в силу того, что число шагов зависит от массива, а сам массив -- от некоторой $\omega$, то и число шагов тоже определяется $\omega$. То есть иначе говоря, число шагов -- это есть некоторая \textit{функция от элементарного исхода} или же, формально, $\xi: \Omega \to \R$.

Однако в классической теории вероятности функция от исхода -- это еще не случайная величина. Но в силу того, что у нас книжка по алгоритмам, а не по основам теории вероятности и задача этого раздела дать хоть какое-то понимание вероятности, мы опустим дополнительные (хоть и очень важные) условия и будем любую функцию от элементарного исхода называть случайной величиной.

Случайные величины делятся глобально на два класса: \textit{дискретные} и \textit{абсолютно непрерывные}. В алгоритмах чаще всего рассматриваются, очевидно, дискретные распределения, хотя и для абсолютное непрерывных иногда находится применение, однако такие случае вылезают максимум на семинарах. Поэтому относительно подробно мы рассмотрим дикретный случай, а абсолютно непрерывный случай остается за рамками нашего курса.

\subsection{Дискретные величины}

Далее считаем что задано некоторое вероятностное пространство $(\Omega, \mathcal{F}, \Pr)$.

\Def Пусть $f: \mathbb{X} \to \mathbb{Y}$ -- произвольное отображение. $\Img f \eqd \{ f(x) \, | \, x \in \mathbb{X} \}$ -- множество значений отображения $f$.

\Def Пусть $\xi: \Omega \to \R$ -- случайная величина. Будем говорить, что $\xi$ -- \textit{дискретная случайная величина}, если она может принимать не более чем счётное количество значений (т.е. множество $\operatorname{Im} \xi$ не более чем счётно). Законом же распределения $\xi$ будем называть такую такую функцию $p: \operatorname{Im} \xi \to [0, 1]$, для которой верно:
\begin{equation*}
    \Pr[\xi = \xi_0] = p(\xi_0), \forall \xi_0 \in \Img \xi
\end{equation*}
И также:
\begin{equation*}
    \sum\limits_{\xi_0 \in \Img \xi} p(\xi_0) = 1
\end{equation*}

Итак, давайте рассмотрим самое популярное распределение в алгоритмах, которое очень желательно знать.

\subsubsection{Равномерное распределение}

Самое наверное известное распределение из всех со школы и его как правило выбирают в качестве модельного распределения при анализе алгоритмов. Такое распределение позволяет смоделировать <<абсолютно случайный непредвзятый выбор элемента>>, т.е. выбор элементов \textit{равновероятно}.

\Def Пусть $\xi: \Omega \to \R$ -- случайная величина и она принимает \textit{конечное} количество значений. Будем говорить, что $\xi$ распределена \textit{равномерно} или же принимает \textit{равновероятно любое значение}, если закон распределения представляется в виде:
\begin{equation*}
    p(\xi_0) = \frac{1}{|\Img \xi|}, \forall \xi_0 \in \Img \xi
\end{equation*}

Например, если рассмотреть бросок кубика и взять за случайную величину число $\eta$, которое выпало, то $\eta$ будет иметь равномерный закон распределения:
\begin{equation*}
    \Pr[\eta = k] = p(k) = \frac{1}{6}, k \in 1 \dots 6
\end{equation*}
Аналогично с монеткой (если принять орла за $0$, а решку за $1$):
\begin{equation*}
    \Pr[\eta = k] = p(k) = \frac{1}{2}, k \in \left\{ 0, 1 \right\}
\end{equation*}

